\begin{helper}
Let \(K\) be a field, let \(p=(x_i)_{i<m}\) be a prefix of vertices of
\(D^*(G(t,K))\), and write \(x_i=(a_i,b_i)\).  Let
\(x=(a,b)\) be a further vertex such that the appended tuple
\(\operatorname{snoc}(p,x)\) is forward independent in \(D^*(G(t,K))\).
Then the representative vector of \(a\) is orthogonal to every vector in
\(W^p(b)\):
\[
  \operatorname{rep}(a)\cdot v=0
  \qquad\text{for every }v\in W^p(b).
\]
\end{helper}
\begin{proof}
By \(\noderef{StarProductForwardIndependentConsistentTuple}\), the forward
independence of the appended tuple is equivalent to
\(D^*\)-consistency of the coordinate tuple in the sense of
\(\noderef{StarProductConsistentTuple}\).  Thus, for every old index \(i<m\),
if \(a_i b\in E(G(t,K))\), then consistency applied to the pair of indices
\(i<m\) and the new last index gives \(a b_i\in E(G(t,K))\).

By \(\noderef{PolarityGraph}\), the edge relation in \(G(t,K)\) is
orthogonality of projective representatives.  Hence the implication just
proved says that every generator \(\operatorname{rep}(b_i)\) occurring in
the spanning set defining \(W^p(b)\) has
\[
  \operatorname{rep}(a)\cdot \operatorname{rep}(b_i)=0.
\]
By \(\noderef{StarProductPrefixSpan}\), \(W^p(b)\) is the \(K\)-linear span
of exactly these generators.  The map
\[
  v\longmapsto \operatorname{rep}(a)\cdot v
\]
is a \(K\)-linear functional, so its zero set is a subspace.  Since all
generators of \(W^p(b)\) lie in this zero set, the whole span \(W^p(b)\)
lies in this zero set.  Therefore
\(\operatorname{rep}(a)\cdot v=0\) for every \(v\in W^p(b)\).
\end{proof}
